\chapter{Metodologia e Plano de Avaliação}\label{chap:metodologia_plano_avaliacao}


% Em Design Science Research, a avaliação é o que sustenta a alegação de utilidade de um artefato e o que distingue a pesquisa de um relato de desenvolvimento \cite{Hevner2004}. O que se avalia em cada agente, como cada medida é obtida e qual critério decide se o resultado sustenta a proposta são fixados antes da execução, para que nenhum resultado seja interpretado depois de conhecido.

% Avalia-se a \textbf{validade interna} de cada agente, isto é, se cada componente faz o que se propõe, e a \textbf{complementaridade} entre eles, isto é, se a análise de conformidade de gênero contribui para a sinalização de risco. Todos os procedimentos são quantitativos, operam sobre evidência interna aos corpora e são executáveis pelo próprio pesquisador, sem recrutamento de participantes. A utilidade percebida por autores em formação, construto perceptual que exige participantes e coleta qualitativa, fica postergada pelas razões expostas na Seção~\ref{sec:utilidade-futura}.


\section{Avaliação do agente de conformidade de gênero}
\label{sec:avaliacao-conformidade}

O agente de conformidade é avaliado em dois níveis, correspondentes às suas camadas. A Especificação de Gênero depende da qualidade da extração de \textbf{moves}, de modo que o segundo nível só é interpretável se o primeiro se sustentar. A Tabela~\ref{tab:niveis-conformidade} resume os dois.

\begin{table}[H]
\centering
\caption{Níveis de avaliação do agente de conformidade de gênero}
\label{tab:niveis-conformidade}
\small
\begin{tabular}{|p{1.0cm}|p{3.4cm}|p{5.4cm}|p{3.5cm}|}
\hline
\textbf{Nível} & \textbf{Construto} & \textbf{Procedimento} & \textbf{Evidência} \\
\hline
1 & Validade da extração de \textbf{moves} (A2) & Comparação com anotação de referência e caracterização dos erros por \textbf{move} & Amostra anotada da trilha-alvo \\
\hline
2 & Validade da Especificação de Gênero (A3--A4) & Validação por retenção, contraste entre trilhas e sensibilidade aos limiares & Trilha-alvo e trilhas de contraste \\
\hline
\end{tabular}
\end{table}

\subsection{Nível 1: validade da extração de \textbf{moves}}
\label{sec:nivel1}

% Toda a cadeia do agente repousa sobre a acurácia da rotulação. Se a extração falha de modo sistemático em um \textbf{move}, as distribuições agregadas em A3 e as regras derivadas em A4 carecem de sentido ainda que sejam computadas corretamente, e o erro se propaga sem deixar rastro nas camadas superiores.

O procedimento consiste na anotação manual de uma amostra de artigos da trilha-alvo segundo o esquema de \textbf{moves} e \textit{steps} adotado, seguida da comparação com os rótulos produzidos pelo modelo. Reportam-se precisão, revocação e medida F1 por \textbf{move}, a medida F1 macro como métrica agregada e a matriz de confusão entre categorias. A matriz importa tanto quanto os índices: confundir \textit{apontar a lacuna} com \textit{ocupar o nicho} tem consequência distinta de deixar de identificar ambos, e apenas a primeira situação preserva a frequência agregada do par.

A execução ocorre em duas etapas. Uma amostra inicial de dez artigos verifica a viabilidade do procedimento e estabiliza o protocolo de anotação, que é fixado antes de qualquer inspeção das saídas do modelo. A avaliação definitiva opera sobre amostra ampliada, estratificada de modo a não concentrar artigos de um mesmo perfil estrutural.

A anotação de referência é produzida por um único anotador, o próprio pesquisador, de modo que não se dispõe de medida de concordância entre anotadores humanos que sirva de teto interpretativo. A alegação sustentável não é que o extrator seja acurado em termos absolutos, mas que concorda com uma anotação de referência em determinada proporção dos casos, com erros concentrados em categorias identificáveis. O produto principal do nível é a caracterização desse padrão de erro e a estimativa de seu efeito sobre as frequências agregadas.

\subsection{Nível 2: validade da Especificação de Gênero}
\label{sec:nivel2}

O segundo nível verifica se a especificação derivada descreve o gênero da trilha ou apenas regularidades acidentais do conjunto processado. Três procedimentos concorrem para a resposta.

\textbf{Validação por retenção.} Parte dos artigos aceitos na trilha-alvo é separada antes da execução do Fluxo A e submetida à verificação do Fluxo B. Se a especificação captura o gênero reconhecido pela comunidade, textos que essa mesma comunidade aceitou devem satisfazer a maior parte das regras de estatuto obrigatório. Uma taxa elevada de reprovação indicaria limiares mal calibrados ou regras sem valor distintivo. A distribuição das ausências observadas entre artigos aceitos informa, ainda, quantas regras um texto reconhecido pela comunidade pode deixar de invocar sem perder pertencimento, o que calibra o critério de suficiência do agente.

Reter artigos reduz o corpus de derivação. Com 67 artigos na trilha-alvo, uma partição de vinte por cento deixa pouco mais de cinquenta para derivar a especificação, e elementos cuja ocorrência esteja próxima de um dos cortes podem mudar de estatuto apenas em função da partição. Compara-se, por isso, a classificação obtida sobre o corpus completo com a obtida sobre o corpus reduzido, e reporta-se toda regra que migre de faixa.

\textbf{Contraste entre trilhas.} A especificação da trilha-alvo é aplicada aos artigos das trilhas de contraste, em dois graus de proximidade. O primeiro confronta a Trilha de Pesquisa em SI com a \textit{Main track} do ENIAC, que pertencem a comunidades e áreas distintas: uma taxa de conformidade substancialmente inferior à dos artigos retidos indica que as regras capturam algo próprio da trilha-alvo, e não propriedades genéricas de artigos publicados em eventos da SBC. O segundo grau é o mais exigente e deriva especificações separadas para a \textit{Main track} e a \textit{Undergraduate track} do ENIAC, comparando-as entre si. As duas compartilham evento, chamada, período, corpo de revisores e área de conhecimento, variando apenas o subgênero. Se as especificações resultantes forem indistinguíveis, a trilha não é o nível de resolução adequado. Reporta-se o conjunto de elementos cujo estatuto difere entre as duas e a magnitude dessa diferença.

\textbf{Sensibilidade aos limiares.} O procedimento percorre a faixa plausível de valores dos dois cortes que convertem frequência em estatuto e reporta quantas e quais regras mudam de classificação, bem como o efeito de cada configuração sobre a validação por retenção. O resultado esperado não é um valor ótimo, mas a identificação de faixas de estabilidade e de pontos de ruptura. Regras que preservam seu estatuto ao longo de toda a faixa constituem o núcleo estável da especificação; regras que oscilam devem ser tratadas com reserva no retorno ao autor.


\section{Avaliação do agente de fronteira de rejeição}
\label{sec:avaliacao-fronteira}

O agente de fronteira dispõe de uma condição que o agente de conformidade não tem: um gabarito real. O ReviewArena registra, para cada artigo, a decisão editorial de aceite ou rejeição, o que permite medir se o retorno de risco discrimina artigos que foram de fato recusados de artigos que foram aceitos. Essa medida responde diretamente se o agente sinaliza algo que importa para o desfecho editorial, e é o eixo desta avaliação. A Tabela~\ref{tab:niveis-fronteira} resume os três níveis.

\begin{table}[H]
\centering
\caption{Níveis de avaliação do agente de fronteira de rejeição}
\label{tab:niveis-fronteira}
\small
\begin{tabular}{|p{1.0cm}|p{3.4cm}|p{5.4cm}|p{3.5cm}|}
\hline
\textbf{Nível} & \textbf{Construto} & \textbf{Procedimento} & \textbf{Evidência} \\
\hline
1 & Validade da extração dos motivos de recusa (Fluxo A) & Comparação com anotação de referência e caracterização dos erros por categoria de motivo & Amostra anotada de avaliações do ReviewArena \\
\hline
2 & Poder discriminante do retorno de risco (Fluxo B) & Classificação aceito/rejeitado sobre partição retida, com métricas de classificação & Partição de teste do ReviewArena \\
\hline
3 & Contribuição da análise de conformidade & Ablação do Fluxo B com e sem a análise de conformidade como insumo & Mesma partição de teste \\
\hline
\end{tabular}
\end{table}

\subsection{Partição do ReviewArena}
\label{sec:particao}

Antes de qualquer execução, o ReviewArena é dividido em duas partições disjuntas. A \textbf{partição de construção} alimenta o Fluxo A e constitui o grafo de conhecimento. A \textbf{partição de teste} é retida e nunca entra no grafo; seus artigos são os manuscritos submetidos ao Fluxo B na avaliação. A separação impede que um artigo de teste seja recuperado como vizinho de si mesmo, o que inflaria artificialmente o resultado, e é a mesma lógica da validação por retenção aplicada ao agente de conformidade.

A partição de teste é estratificada pela decisão editorial, de modo a conter proporções conhecidas de artigos aceitos e rejeitados, e pelo veículo, para que a avaliação não se concentre em uma única conferência da base. A proporção entre as partições é fixada antes da execução e reportada junto aos resultados.

\subsection{Nível 1: validade da extração dos motivos de recusa}
\label{sec:nivel1-fronteira}

O grafo de conhecimento é tão bom quanto os motivos de recusa que o Fluxo A extrai das avaliações. Se a extração confunde categorias de motivo ou deixa de identificar motivos presentes, os vizinhos recuperados carregarão rótulos errados, e o retorno ao autor reportará riscos que os revisores não apontaram.

O procedimento espelha o Nível 1 do agente de conformidade. Uma amostra de avaliações da partição de construção é anotada manualmente segundo um esquema de categorias de motivo de recusa, fixado antes da inspeção das saídas do modelo. Os motivos extraídos pelo modelo são comparados à anotação, reportando-se precisão, revocação e medida F1 por categoria, a medida F1 macro como agregado e a matriz de confusão. Como no outro agente, a anotação é produzida por um único anotador, e a alegação sustentável é de concordância com uma referência, não de acurácia absoluta.

\subsection{Nível 2: poder discriminante do retorno de risco}
\label{sec:nivel2-fronteira}

Cada artigo da partição de teste é submetido ao Fluxo B como se fosse um manuscrito. O agente produz o retorno de risco, e desse retorno deriva-se uma predição binária, sinalizado como exposto a risco de recusa ou não, por regra fixada antes da execução e aplicada de modo uniforme a todos os artigos. A predição é comparada à decisão editorial real registrada na base.

Trata-se de um problema de classificação binária, e reportam-se as métricas usuais: precisão, revocação e medida F1 para a classe de rejeitados, a acurácia global e a matriz de confusão. A medida F1 da classe de rejeitados é a métrica principal, porque é a que expressa se o agente sinaliza os artigos que de fato foram recusados sem sinalizar indistintamente todos os artigos. Como as classes podem ser desbalanceadas, reporta-se também a área sob a curva ROC (Receiver Operating Characteristic), obtida a partir do grau de risco que o agente atribui, e o desempenho de uma linha de base trivial que sinalize todos os artigos como expostos, para que o ganho do agente seja lido contra ela.

O nível responde à pergunta central da avaliação: se o retorno de risco distingue artigos aceitos de rejeitados acima do acaso e acima da linha de base trivial, o agente sinaliza algo que corresponde ao desfecho editorial real.

% \subsection{Nível 3: contribuição da análise de conformidade}
% \label{sec:nivel3-fronteira}

% A proposta afirma que os dois agentes se complementam, e que a análise de conformidade de gênero entra no Fluxo B do agente de fronteira como sinal de aproximação. Esse nível verifica se essa entrada contribui de fato, por meio de uma ablação controlada.

% O Nível 2 é executado duas vezes sobre a mesma partição de teste, com o mesmo grafo, o mesmo modelo e a mesma regra de predição, variando apenas o insumo: em uma condição o Fluxo B recebe o manuscrito e a análise de conformidade; na outra, recebe apenas o manuscrito. Comparam-se as métricas de classificação das duas condições. Se a condição com a análise de conformidade apresentar medida F1 superior para a classe de rejeitados, a complementaridade está demonstrada; se as duas condições forem indistinguíveis, a análise de conformidade não contribui para a sinalização de risco, e a alegação de complementaridade não se sustenta.

% A ablação é controlada porque as duas condições diferem em um único fator e são medidas contra o mesmo gabarito, o que permite atribuir a diferença observada à presença da análise de conformidade.


\section{Critérios de sucesso}
\label{sec:criterios}

%Cada nível produz números, e números só decidem alguma coisa se o critério estiver fixado antes da execução. 
A Tabela~\ref{tab:criterios} registra os critérios adotados. Eles são decisões de projeto, e não constantes derivadas da literatura; sua função é impedir que o resultado seja interpretado depois de conhecido. Os patamares numéricos são fixados nos protocolos de anotação e de partição, antes da execução, e reportados junto aos resultados.

\begin{table}[H]
\centering
\caption{Critérios de sucesso por agente e nível}
\label{tab:criterios}
\small
\begin{tabular}{|p{2.2cm}|p{5.0cm}|p{6.2cm}|}
\hline
\textbf{Agente, nível} & \textbf{Critério} & \textbf{Interpretação de um resultado adverso} \\
\hline
Conformidade, 1 & Medida F1 macro acima do patamar fixado, sem \textbf{move} obrigatório abaixo de um piso individual & Extração inadequada para sustentar a especificação; exige revisão do esquema de rotulação ou do modelo antes de prosseguir \\
\hline
Conformidade, 2 & Maioria dos artigos retidos satisfaz as regras obrigatórias; conformidade das trilhas de contraste inferior à dos retidos; núcleo de regras estável na faixa de limiares & Especificação sem poder discriminante, descrevendo artigos científicos em geral e não o gênero da trilha \\
\hline
Fronteira, 1 & Medida F1 macro acima do patamar fixado para as categorias de motivo & Extração inadequada; os motivos do grafo não refletem o que os revisores apontaram \\
\hline
Fronteira, 2 & Medida F1 da classe de rejeitados e área sob a curva ROC acima da linha de base trivial & O retorno de risco não distingue aceitos de rejeitados; o agente não sinaliza o que importa para o desfecho editorial \\
\hline
% Fronteira, 3 & Condição com análise de conformidade supera a condição sem análise na medida F1 da classe de rejeitados & A análise de conformidade não contribui para a sinalização de risco; a alegação de complementaridade não se sustenta \\
% \hline
\end{tabular}
\end{table}

% Um resultado adverso em qualquer nível é achado da pesquisa, e não falha a ser contornada. Estabelecer que a modelagem de gênero por derivação empírica não sustenta a verificação de conformidade, ou que avaliações negativas de outras comunidades não permitem sinalizar risco para uma trilha brasileira, seria contribuição legítima ao conhecimento, ainda que negativa.


\section{Avaliação de utilidade percebida}
\label{sec:utilidade-futura}

A utilidade percebida do artefato por autores em início de formação é construto perceptual e não admite substituto computacional: requer participantes, seus próprios manuscritos e coleta de dados qualitativos. Está prevista para a etapa subsequente do trabalho, como estudo exploratório com um pequeno número de mestrandos do programa, mediante termo de consentimento livre e esclarecido e observância dos procedimentos institucionais aplicáveis.

A postergação é escolha de escopo. Os níveis aqui descritos estabelecem que cada agente mede o que se propõe a medir e que a análise de conformidade contribui para a sinalização de risco. Avaliar utilidade antes dessas condições produziria percepções sobre um sistema em construção, e não sobre o artefato final. 


% \section{Ameaças à validade}
% \label{sec:ameacas}

% % \textbf{Anotador único.} As anotações de referência dos dois agentes são produzidas pelo pesquisador, o que implica ausência de medida de concordância e risco de viés na direção do comportamento esperado do extrator. Mitiga-se fixando os protocolos de anotação antes da inspeção dos resultados do modelo e reportando o padrão de erro em lugar de acurácia absoluta.

% % \textbf{Ausência de validação por especialistas.} A pertinência das regras de gênero não é submetida a revisores experientes da trilha, o que impede detectar regras que a comunidade reconheça e o corpus não revele. O Nível 2 do agente de conformidade atesta que as regras descrevem o corpus e discriminam entre trilhas, não que sejam exaustivas.

% % \textbf{Viés de sobrevivência no agente de conformidade.} O corpus desse agente contém apenas artigos aceitos. A validação por retenção estabelece que a especificação é compatível com textos aprovados, não que a não conformidade implique rejeição. A fronteira entre aceitação e recusa é medida pelo agente de fronteira, e apenas por ele.

% \textbf{Deslocamento de domínio no agente de fronteira.} O gabarito do agente de fronteira provém do ReviewArena, cujas avaliações são de conferências internacionais de aprendizado de máquina e processamento de linguagem natural. O poder discriminante medido no Nível 2 refere-se a essas comunidades. Nenhum resultado deste capítulo autoriza afirmar que o agente discrimina aceitos de rejeitados em uma trilha brasileira, o que exigiria um corpus de avaliações brasileiras. A avaliação demonstra o mecanismo, não sua transferência ao contexto-alvo.

% \textbf{Regra de predição.} A conversão do retorno de risco em predição binária depende de uma regra fixada pelo pesquisador. Regras distintas produzem métricas distintas. Mitiga-se fixando a regra antes da execução e reportando a sensibilidade das métricas a variações dessa regra, do mesmo modo que se reporta a sensibilidade aos limiares no agente de conformidade.

% \textbf{Dependência entre níveis.} Erros sistemáticos de extração, em qualquer dos agentes, propagam-se às camadas superiores e podem produzir resultados internamente consistentes e substantivamente equivocados. A caracterização do erro obtida no Nível 1 de cada agente deve acompanhar a leitura dos demais níveis.

% \textbf{Uma edição por trilha.} As especificações de gênero derivam de uma única edição de cada evento. Regularidades próprias daquela edição não se distinguem de regularidades estáveis da trilha, e o contraste entre trilhas atenua apenas parcialmente a limitação, uma vez que ambas as trilhas do ENIAC provêm da mesma edição.

% \textbf{Alcance das alegações.} Os procedimentos operam sobre uma trilha-alvo, uma língua e uma comunidade, com o agente de fronteira avaliado em comunidades distintas da alvo. Não se sustenta alegação de generalidade a outros gêneros acadêmicos, a outras comunidades científicas ou a outros idiomas.


% \section{Síntese}
% \label{sec:sintese-avaliacao}

% A avaliação responde a quatro perguntas encadeadas. Para o agente de conformidade, se ele enxerga no texto o que afirma enxergar e se o que deriva do corpus descreve o gênero de uma trilha em particular. Para o agente de fronteira, se os motivos de recusa que extrai correspondem aos que os revisores apontaram, se o retorno de risco distingue artigos aceitos de rejeitados, e se a análise de conformidade contribui para essa distinção. Respondidas afirmativamente, estabelecem a validade interna de cada agente e a complementaridade entre eles, condição para que a avaliação de utilidade com autores produza percepções atribuíveis a alguma camada do sistema. Respondidas negativamente, indicam qual camada revisar antes de prosseguir.